% （本节重点见 §2 开头，由分片 C 负责）

\subsection{规范场的 $N=2$ 超对称拉格朗日量}

壳上的 $N=1$ 标量多重态 $(A, \psi)$ 与矢量多重态 $(A_\mu,\lambda)$
合在一起，其场内容与壳上的 $N=2$ 矢量多重态 $(A,\psi,\lambda, A_\mu)$
相同。拉格朗日量 (\ref{fullC}) 包含所有这些场，但本身并不具有
$N=2$ 超对称性。现在假定 $(A,\psi,\lambda, A_\mu)$ 构成一个
$N=2$ 矢量多重态，并讨论这一假定对 (\ref{fullC}) 中的
$N=1$ 拉格朗日量所施加的限制。首先，既然 $A^a_\mu$ 与 $\lambda^a$
属于规范群的伴随表示，那么 $A_i$ 与 $\psi_i$ 若要成为同一多重态的
一部分，也必须属于同一表示。因此 $T^a_{ij}=-if^a_{ij}$，并且指标集合
$\{ i \}$ 与 $\{ a \}$ 重合。其次，既然 $N=2$ 代数中的两个超对称生成元
地位相同，(\ref{fullC}) 中的费米子 $\psi^a$ 与 $\lambda^a$ 也应当如此。
为满足这一条件，我们令超势 ${\cal W}$ 为零，因为它只与 $\psi^a$ 耦合。
这一条件同时固定了拉格朗日量中 Yang--Mills 部分与标量部分之间任意的
相对归一化，因为它要求两个费米子的动能项具有相同的归一化。这可以通过
在 (\ref{fullC}) 中作标度变换 $\Phi\rightarrow\Phi/g$ 来实现。结果表明，
只要拉格朗日量 (\ref{fullC}) 满足这些条件，它就具有 $N=2$ 超对称性。
此时含辅助场的各项取形式
$$
\frac{1}{g^2}{\rm Tr}\left(\frac{1}{2}DD +D\,[A^\dagger\,,\,A] + 
F^\dagger F \right)\,, 
$$
其中我们用了记号 $\Phi=\Phi^a T^a$，$T^a$ 取基本表示。消去 $D$ 与 $F$ 后，
得到标量势
\begin{equation}
V=- \frac{1}{2g^2} {\rm Tr}\left([A^\dagger\,,\,A]^2 \right)\,.
\label{scalarV}
\end{equation}
% [Tutor Note: 见原文 PDF 第 33 页式 (71)。这里的符号与因子极易出错。把辅助场
% 部分的拉格朗日量写全为 $\frac{1}{g^2}{\rm Tr}(\frac{1}{2}D^aD^a+D^a[A^\dagger,A]^a+F^{\dagger a}F^a)$，
% $D$ 的代数方程给出 $D^a=-[A^\dagger,A]^a$；代回后 $\frac{1}{2}D^2$ 项只剩
% $-\frac{1}{2g^2}{\rm Tr}([A^\dagger,A]^2)$。由于 $[A^\dagger,A]$ 反厄米，这个量非负，
% 这正是标量势的正定性来源。注意：(i) 系数 $1/2$ 与 $1/g^2$ 同时出现，漏掉任何一个都会
% 使 (71) 与 (72) 中的 $-\frac{1}{2}[A^\dagger,A]^2$ 对不上；(ii) 原文把二次项写成
% $\frac{1}{2}DD$（$D$ 的指标缩并省略），不要误读为两个不同的 $D$ 场。]

具有 $N=2$ 超对称性的完整拉格朗日量现在可以写成
\begin{eqnarray}
{\cal L}&=&\frac{1}{8\pi}{\rm Im\, Tr}\left[\,\tau\left(\int d^2\theta\,
W^\alpha W_\alpha +2\,\int d^2\theta d^2\bar\theta\, \Phi^\dagger 
e^{-2V}\Phi \right)\right] 
\nonumber \\ 
&=&\frac{1}{g^2}{\rm Tr}\Big(-\frac{1}{4}F_{\mu\nu}F^{\mu\nu}+ 
g^2\frac{\theta}{32\pi^2}F_{\mu\nu}\widetilde F^{\mu\nu}+
(D_\mu A)^\dagger D^\mu A -\frac{1}{2} [A^\dagger\,,\,A]^2
\nonumber\\
&&{\hphantom {\frac{1}{g^2}{\rm Tr}(\,} }
-i\,\lambda\sigma^\mu D_\mu\bar\lambda -i\,\bar{\psi} 
\bar{\sigma}^\mu D_\mu\psi-i\sqrt{2}\,[\lambda,\psi]\,A^\dagger
-i\sqrt{2}\,[\bar\lambda,\bar\psi]\,A \Big)\,,
\label{N=2}
\end{eqnarray}
其中在分量展开中已消去辅助场。
% [Tutor Note: 见原文 PDF 第 33 页式 (72)。(72) 的第一行是全纯规范耦合
% $\tau=\theta/2\pi+4\pi i/g^2$ 的 $N=1$ 超空间写法；第二行才是分量形式，其 Yang--Mills
% 项系数为 $1/g^2$，标量动能项系数也是 $1/g^2$，两者的相对归一化正是上一段所说的
% "两个费米子动能归一化相同" 所固定。与 §1 的 Yang--Mills--Higgs 拉格朗日量逐项对照时，
% 注意这里 $A$ 带有因子 $1/g$（见正文对 $\Phi\to\Phi/g$ 的标度），所以 Higgs 势写成
% $-(1/2g^2){\rm Tr}([A^\dagger,A]^2)$ 而非 §1 中的 $\lambda(\phi^\dagger\phi-v^2)^2$ 形式。
% 另外 $\theta$ 项系数 $g^2\theta/32\pi^2$ 中的 $g^2$ 来自 $\tau$ 的虚部与 $F\widetilde F$
% 的归一化，切勿与 $\theta/32\pi^2$ 混淆。]
注意 (\ref{N=2}) 是 §1 所述 Yang--Mills--Higgs 理论的超对称推广：Higgs 势 $V$
由 (\ref{scalarV}) 给出，并且带有一个 $\theta$ 项。因此 §1 的讨论同样适用于我们的
$N=2$ 超对称理论。特别地，Higgs 真空由 $D_\mu A=0, V=0$ 定义，并且只要 $A$ 与
$A^\dagger$ 对易，势在非零场位形上就为零。于是该模型既有磁单极解也有双子解，并包含
有质量的规范玻色子。例如，若规范群为 $SU(2)$ 或 $SO(3)$，它破缺到 $U(1)$，
其中两个规范玻色子获得质量。

现在假定我们关心的是该模型在低于某个截断 $\Lambda$ 的能量下的行为，而 $\Lambda$
小于理论中最轻的有质量态的质量。在这样的能量下不会遇到任何壳上有质量态，物理可由
威尔逊低能有效作用量描述。该有效作用量通过把全部有质量态以及标度 $\Lambda$ 以上的
全部无质量激发完全积分掉而得到 \cite{SV1,SV2}。这通常是一个复杂的过程，无法显式
完成。所幸在我们这个模型中，威尔逊有效作用量的一般形式受到 $N=2$ 超对称性的严格
限制。把理论 (\ref{N=2}) 用下面描述的 $N=2$ 超场表述时，这一点最容易看出。
 
\noindent\underline{$N=2$ 超空间表述：}
$N=2$ 超空间是通过在 $N=1$ 超空间上再加四个费米自由度（例如 $\tilde\theta$ 与
$\bar{\tilde\theta}$）而得到的。因此，一般的 $N=2$ 超场可以写成
$F(x,\theta,\bar\theta,\tilde\theta,\bar{\tilde\theta})$。
我们需要一个分量与 $N=2$ 矢量多重态相同的超场。这可通过对一般 $N=2$ 超场施加
手征性约束与实性约束而得到 \cite{Grimm}。下面我们按 \cite{bilal} 的方式来简要
描述这一点。$N=2$ 手征超场 $\Psi$ 由约束 $\bar
D_{\dot\alpha}\Psi=0$ 与 $\bar{\tilde D}_{\dot\alpha}\Psi=0$ 定义，其中
超协变导数 $\bar{\tilde D}_{\dot\alpha}$ 的定义与 $\bar D_{\dot\alpha}$ 相同，
只需把 $\theta$ 换成 $\tilde\theta$。$\Psi$ 的展开可以按 $\tilde\theta$ 的幂次
整理，并写成
$$
\Psi=\Psi^{(1)}(\tilde y,\theta)+\sqrt{2}\tilde\theta^\alpha
\Psi^{(2)}_\alpha(\tilde y, \theta) + \tilde\theta^\alpha
\tilde\theta_\alpha \Psi^{(3)}(\tilde y, \theta)\,,
$$
其中 $y^\mu =x^\mu +i\theta\sigma^\mu\bar\theta +i\tilde\theta
\sigma^\mu\bar{\tilde\theta}$。这一展开有助于把 $N=2$ 表述与我们迄今一直使用的
$N=1$ 语言联系起来。显然，分量 $\Psi^{(1)}$ 与 $N=1$ 手征超场 $\Phi$ 形式相同。
其余两个分量受一个实性条件的约束。结果是 $\Psi^{(2)}_\alpha=
W_\alpha (\tilde y, \theta)$，即 (\ref{W-nonabelian}) 所给，而 $\Psi^{(3)}$ 由
下式给出
$$
\Psi^{(3)}(\tilde y,\theta)= \Phi^\dagger(\tilde y-i\theta\sigma 
\bar\theta, \theta,\bar\theta) exp \left[2gV(\tilde y -i\theta\sigma
\bar\theta, \theta,\bar\theta\right]\,\mid_{\bar\theta\bar\theta}\,.
$$
% [Tutor Note: 见原文 PDF 第 34 页 $\Psi^{(3)}$ 的表达式。原文此处排版有笔误：
% $\exp$ 的方括号内 $2gV(\cdots$ 只开了一个圆括号而没有闭合，源码中
% $\backslash$left[ $\cdots$ $\backslash$right] 也不配对；按上下文（以及 $\bar\theta\bar\theta$
% 投影）应为 $\Phi^\dagger(\cdots)\exp[2gV(\cdots)]|_{\bar\theta\bar\theta}$。
% 本书译文照录原文的写法，未作修补。]
这里 $\Phi(\tilde y -i\theta\sigma\bar\theta,\theta,\bar\theta)$ 应按
(\ref{csfield}) 中的展开来理解。显然，$\Psi$ 具有与 $N=2$ 矢量多重态相同的场内容。
可以验证，用 $N=2$ 超场 $\Psi$ 表示时，$N=2$ 拉格朗日量 (\ref{N=2}) 由如下
紧凑表达式给出
\begin{equation}
{\cal L}= \frac{1}{4\pi}{\rm Im\,Tr}\int d^2\theta d^2\tilde\theta\,
\frac{1}{2}\tau\Psi^2\,.
\label{N=2manifest}
\end{equation}

利用 $N=2$ 手征超场 $\Psi$（也称为 $N=2$ 矢量超场），我们现在可以构造规范场的
最一般的 $N=2$ 拉格朗日量：对应于任意函数 ${\cal F}(\Psi)$，都可以构造一个
拉格朗日量
\begin{eqnarray}
{\cal L}&=& \frac{1}{4\pi}{\rm Im\,Tr}\int d^2\theta d^2\tilde\theta
\,{\cal F}(\Psi) 
\nonumber\\
&=&\frac{1}{8\pi}{\rm Im}\left(\int d^2\theta\, {\cal F}_{ab}(\Phi) 
W^{a\alpha}W^b_\alpha + 2\int d^2\theta d^2\bar\theta\, (\Phi^\dagger 
e^{2gV})^a {\cal F}_a(\Phi)\right)\,. 
\label{N=2general}
\end{eqnarray}
这里 ${\cal F}_a(\Phi)=\partial{\cal F}/\partial\Phi^a$，
${\cal F}_{ab}(\Phi)=\partial^2{\cal F}/\partial\Phi^a\partial\Phi^b$，
${\cal F}$ 称为 $N=2$ 预势。由上可以直接读出凯勒势为
${\rm Im}(\Phi^{\dagger a} {\cal F}_a(\Phi)$。它给出场空间上的度规
$g_{ab}={\rm Im}(\partial_a\partial_b{\cal F})$。这种形式的度规称为特殊凯勒度规。
% [Tutor Note: 见原文 PDF 第 34 页式 (74) 及其后。两处需留意。
% (i) 原文 (74) 第二行写的是 $(\Phi^\dagger e^{2gV})^a{\cal F}_a(\Phi)$，指数是
% $+2gV$；而 (72) 中规范不变的凯勒项是 $\Phi^\dagger e^{-2V}\Phi$。若取
% ${\cal F}=\frac12\tau\Psi^2$ 回到 (72)，按规范不变性这里通常应是 $\Phi^\dagger e^{-2gV}$
% （$g$ 因 $\Phi\to\Phi/g$ 的标度而出现；同一理由下 (72) 的指数也应是 $-2gV$）。
% 也就是说 (72) 与 (74) 两处的符号及 $g$ 因子互不自洽，至少有一处是原文笔误。
% 译文照录原文，读者自行推导有效作用量时请以规范不变性为准核对文献 [27,28]
% （Grimm；Bilal）的约定。
% (ii) 紧接着"读出凯勒势"的句子在原文中括号不配对：${\rm Im}(\Phi^{\dagger a}{\cal F}_a(\Phi)$
% 少一个右括号，译文照录；另外度规 $g_{ab}={\rm Im}(\partial_a\partial_b{\cal F})$
% 严格说是混合分量 $g_{a\bar b}$ 的写法（凯勒度规的非零分量是 $\partial_a\partial_{\bar b}K$），
% 原文指标约定如此。]
如果我们要求可重整性，那么 ${\cal F}$ 必须如 (\ref{N=2manifest}) 那样是 $\Psi$
的二次式。然而，如果要写出低能有效作用量，可重整性就不再是判据，${\cal F}$ 可以有
更复杂的形式。特别地，我们可以从对应于二次预势的微观理论 (\ref{N=2}) 出发，尝试为
低能威尔逊有效作用量构造修改后的 ${\cal F}$。这一函数的精确确定正是 Seiberg 与
Witten 工作的主题。
% [Tutor Note: 见原文 PDF 第 34 页末段。这里的逻辑是：$N=2$ 超对称把低能有效作用量
% 完全装入一个全纯函数 ${\cal F}(\Psi)$（预势）之中，因此确定量子修正后的 ${\cal F}$
% 就等于确定整个低能理论。注意 ${\cal F}$ 与 $N=1$ 的超势 ${\cal W}$ 是两个不同的
% 对象（前者是全纯规范动力学的一部分，后者给出 Yukawa 耦合与 F 项），全文中
% "prepotential" 一律译作"预势"。可重整性只允许二次 ${\cal F}$，故低能有效理论
% 中 ${\cal F}$ 的非二次部分必定来自量子修正，确定它便是非微扰问题。]

\subsection{物质场的 $N=2$ 超对称拉格朗日量}

由于物质场与规范场在规范群的不同表示下变换，它们不能属于同一多重态。
$N=2$ 的物质超多重态称为超多重态，它包含一对复标量与一对二分量旋量，
它们都在规范群的同一表示下变换。由前面关于超对称代数表示的讨论可知，
若超多重态有质量，则其质量应当以超对称代数中中心扩张的形式出现。
% [Tutor Note: 见原文 PDF 第 34 页 §2.7 开头。这里的论证是 §2.3 关于含中心荷
% 代数表示（原文 PDF 第 26 页起）的结论的直接应用：无中心荷时，有质量的
% $N=2$ 多重态必须是"长"多重态，其分量数比超多重态（短多重态）多；只有中心荷
% 取特定值（BPS 条件 $M=\sqrt{2}|Z|$）时，有质量表示才能保持"短"，才能装进
% 超多重态。因此裸质量 $m_i$ 必须同时作为中心荷出现——这正是 §2.9 要计算的内容。]
这是因为否则 $N=2$ 超对称性将要求比超多重态所含更多的分量。用 $N=1$ 记号，
一个超多重态包含一个手征超场 $Q$ 和一个反手征超场
$\widetilde Q^\dagger$，二者都在规范群 $SU(N_c)$ 的同一表示 $N_c$ 下变换。
我们把 $Q$ 与 $\widetilde Q$ 的分量分别记为 $(q,\psi_q,F_q)$ 与
$(\widetilde q,\psi_{\widetilde q},F_{\widetilde q})$。$N_f$ 个超多重态
（以指标 $i$ 标记）与一个 $N=2$ 矢量多重态相互作用的拉格朗日量，可以部分地
从 $N=1$ 理论推断出来，其形式为
\begin{eqnarray}
{\cal L}=&&\int d\theta^4\left( Q^\dagger_i e^{-2V}Q_i +\widetilde Q_i
e^{2V}\widetilde Q^\dagger_i\right)+\int
d\theta^2\left(\sqrt{2}\widetilde Q_i\Phi Q_i+m_i\widetilde Q_i
Q_i\right)+\,h.c.\nonumber\\       &&+\cdots
\label{N=2hyper}
\end{eqnarray}
这里省略号代表纯 $N=2$ 矢量多重态的拉格朗日量，并且我们略去了规范群指标。
项 $\sqrt{2}\widetilde Q_i\Phi Q_i$ 通过 $N=2$ 超对称性与超多重态同
$N=1$ 矢量多重态 $V$ 的耦合相联系。在 $N=1$ 拉格朗日量中并不要求存在这一项。
与通常的 QCD 一样，当所有质量 $m_i$ 都相等时，理论在整体味群 $SU(N_f)$ 下不变。
% [Tutor Note: 见原文 PDF 第 35 页式 (75)。这是本片最易写错的超势之一，请注意：
% (i) 超势为 $\sqrt{2}\widetilde Q_i\Phi Q_i+m_i\widetilde Q_iQ_i$，其中 $\Phi=\Phi^aT^a$
% 取伴随表示，$Q_i$ 与 $\widetilde Q_i$ 分别取 $N_c$ 与 $\bar N_c$ 表示，指标缩并后
% 是规范单态。(ii) 系数 $\sqrt2$ 由 $N=2$ 超对称固定，不能与质量项的相对系数任意
% 改变；它与式 (78) 中物质超流里的 $i\sqrt2$ 因子同源。(iii) 当 $A$ 取真空期望值
% $a$ 时，$\sqrt2\widetilde q\,\Phi q$ 与 $m\widetilde q q$ 一起给出 $N=2$ 不变
% 的质量项，其结构保证了 BPS 关系 $M=\sqrt2|Z|$（见 §2.9）。(iv) 原文把 $N_f$ 个
% 超多重态的质量都取为对角形式，未写味混合；一般情形可通过对 $Q_i$ 作 $U(N_f)$ 旋转
% 对角化。]

消去超多重态中出现的辅助场 $F_q$ 与 $F_{\widetilde q}$，会给出对标量势的一项
贡献（与 (\ref{Spot}) 比较）
$$
V=\frac{1}{2}g^2\sum_a D_a D^a\,,
$$
其中
\begin{equation}
D^a=\sum_{i=1}^{N_f}\left(q^\dagger_i\lambda^a q_i -\widetilde
q_i\lambda^a \widetilde q_i^\dagger\right)\,.
\label{D-term}
\end{equation}
这里 $\lambda^a$ 是规范群在基本表示下的生成元。这一项称为 $D$ 项。需要指出的是，
$N=2$ 代数具有整体 $SU(2)_R$ 对称性，它理应是拉格朗日量的对称性。然而，在我们
用 $Q$ 与 $\widetilde Q^\dagger$ 对超多重态所作的分解中，这一对称性并不显然，
因为它把标量分量 $q$ 与 $\widetilde q^\dagger$ 作为二重态旋转。写出具有显式
$SU(2)_R$ 对称性的 $N=2$ 拉格朗日量并不困难 \cite{West}。
% [Tutor Note: 见原文 PDF 第 35 页式 (76)。两处易错：(i) $D^a$ 中 $q^\dagger_i\lambda^a q_i$
% 与 $-\widetilde q_i\lambda^a\widetilde q_i^\dagger$ 两项符号相反，这是因为 $Q$ 与
% $\widetilde Q$ 的 $U(1)$ 荷相反；若把第二项的符号写错，消去 $D$ 后的标量势就不再
% 半正定。(ii) 这里的 $\lambda^a$ 是基本表示的生成元矩阵（$N_c\times N_c$），
% 不是伴随表示的结构常数 $f^{abc}$；$\sum_a$ 是对规范群所有生成元求和，
% 与式 (71) 中 ${\rm Tr}$ 的写法只是记号不同。(iii) 原文把 $D$ 项与超多重态质量项
% 一起极小化并引用文献 [29,30]；这一步的物理内容是寻找标量势的平坦方向，
% 其分析结果直接给出式 (77) 的矩阵条件。]

\noindent\underline{$D$ 项的极小化：}
为求理论的真空，现在还必须极小化 $D$ 项贡献与超多重态质量项对标量势的贡献
\cite{DDS,ADS}。对于非零的夸克质量，唯一的解是
$q=\widetilde q =0$。因此只有标量 $A$ 可以有非零的真空期望值。但当 $m_i=0$ 时，
$D$ 项可以在一组非零的 $q$ 与 $\widetilde q$ 处取极小（势有平坦方向）。然而，
拉格朗日量中的 $\widetilde Q\Phi Q$ 项此时要求 $A_{vac}=0$。下面我们确定这些
平坦方向。

$D$ 项的极小值对应于 $D^a=0$。为求解此式，注意 $q^{(i)}_\alpha$（其中 $\alpha$
表示色指标）可以看作 ${\bf C}^{N_f}$ 中的 $N_c$ 个矢量。因此可以构造标量积矩阵
$$
\sum_i q^{(i)}_\alpha q^{\dagger (i)}_\beta=q_\alpha\cdot
q^\dagger_\beta=(qq^\dagger)_{\alpha\beta}\,.
$$
用它表示，
$$
\sum_i q^\dagger_i\lambda^aq_i={\rm Tr}\,\left(qq^\dagger
\lambda^a\right) \,.
$$
类似地，我们构造矩阵 $\widetilde q_\alpha^\dagger \cdot
\widetilde q_\beta = (\widetilde q^\dagger\widetilde q)_{\alpha\beta}$。
于是 $D^a=0$ 可以写成
$$
D^a={\rm Tr}\left[\left(qq^\dagger-\widetilde q^\dagger
\widetilde q\right)\,\lambda^a\right]=0 \,.
$$
由于 $\lambda^a$ 处于 $SU(N_c)$ 的不可约表示，上述条件由下式解出
\begin{equation}
(qq^\dagger) - (\widetilde q^\dagger\widetilde q)=c^2{\bf 1}_{N_c}\,.
\label{D=0} 
\end{equation}
这里可以区分两种情形：$N_f < N_c$ 与 $N_f\geq N_c$。
% [Tutor Note: 见原文 PDF 第 36 页式 (77) 之前。由 ${\rm Tr}[(qq^\dagger-\widetilde q^\dagger\widetilde q)\lambda^a]=0$
% 推出 (77)，用的是 $\lambda^a$ 在不可约表示下的完备性（Schur 引理）：若矩阵 $M$ 与
% 所有生成元对易，则 $M\propto{\bf 1}$。这里 $M=qq^\dagger-\widetilde q^\dagger\widetilde q$
% 实际上满足的是更强的条件 ${\rm Tr}(M\lambda^a)=0$（对所有 $a$），即 $M$ 正交于
% 所有无迹生成元，故 $M\propto{\bf 1}$。注意常数写成 $c^2{\bf 1}_{N_c}$（$c$ 可正可负，
% 写 $c^2$ 便于下面开方），这种"用平方记常数"的写法在原文中出现多次。]

\noindent 1) $N_f < N_c$：此时矩阵 $(qq^\dagger)$ 的秩为 $N_f$，因而它有
$N-N_f$ 个零本征值。对矩阵 $(\widetilde q^\dagger\widetilde q)$ 也一样。利用
$SU(N_c)\times SU(N_f)\times U(1)_R$ 旋转，可以把 $(qq^\dagger)$ 对角化，而
条件 (\ref{D=0}) 于是要求 $(\widetilde
q^\dagger\widetilde q)$ 也必须是对角的。因此在这个基下，可以解出 $q$ 与
$\widetilde q$ 为
$$
q = \left(
\begin{array}{cccc}
v^{(1)}_1 & 0 &\cdots&0\\
0 & v^{(2)}_2 &&\\
\vdots&&\ddots & \\
&&&v^{(N_f)}_{N_f}\\
0&0&\cdots&0\\
\vdots&\vdots&&\vdots
\end{array}
\right)\,, \quad
\widetilde q = \left(
\begin{array}{cccccc}
\tilde v^{(1)}_1 & 0 & \cdots & & 0 & \cdots \\
 0 & \tilde v^{(2)}_2 & & & 0 & \cdots \\
\vdots  & & \ddots &  &   &  \\
 0 & & &\tilde v^{(N_f)}_{N_f} & 0 & \cdots  \\
\end{array}
\right)\,.
$$
零本征值的存在意味着 $c = 0$，从而
$v^{(i)}_i =\tilde{v}^{(i)}_i$。此时规范对称性破缺到 $SU(N_c-N_f)$，
但 $N_c=N_f-1$ 时除外——那时它完全破缺。$2N_fN_c-N^2_f$ 个夸克超场变重，
剩下的 $N^2_f$ 个夸克超场保持无质量，对应于破缺的整体对称性的 Goldstone 玻色子。
% [Tutor Note: 见原文 PDF 第 36 页第 1 种情形。此段有两处原文笔误，译文照录：
% (i) "它有 $N-N_f$ 个零本征值"应为 $N_c-N_f$（同段其余各处均用 $N_c$）；
% (ii) "但 $N_c=N_f-1$ 时除外，那时它完全破缺"应为 $N_c=N_f+1$：因为破缺后的
% 未破缺群是 $SU(N_c-N_f)$，只有当 $N_c-N_f=1$ 时它是平凡群 $SU(1)$，规范对称性
% 才完全破缺。另外，"$2N_fN_c-N_f^2$ 个夸克超场变重"这个数目恰等于被破缺的规范
% 生成元数 $\dim SU(N_c)-\dim SU(N_c-N_f)$，即 Higgs 机制中每个被吃掉一个手征超场；
% 这与"剩下 $N_f^2$ 个无质量超场作为破缺整体对称性的 Goldstone 玻色子"是自洽的。]

\noindent 2) $N_f \leq N$:  此时 $(qq^\dagger)$
的秩为 $N_c$，一般说来它的本征值不为零。于是 $c\not=0$。与上面同样的论证，
此时 $q$ 与 $\widetilde q$ 的真空值可以写成
$$
q = \left(
\begin{array}{cccccc}
v^{(1)}_1 & 0 & \cdots & & 0 & \cdots \\
 0 & v^{(2)}_2 & & & 0 & \cdots \\
\vdots  & & \ddots &  &   &  \\
 0 & & & v^{(N_c)}_{N_c} & 0 & \cdots  \\
\end{array}
\right)\,,\quad
\widetilde q = \left(
\begin{array}{cccc}
\tilde v^{(1)}_1 & 0 &\cdots&0\\
0 &\tilde v^{(2)}_2 &&\\
\vdots&&\ddots & \\
&&&\tilde v^{(N_c)}_{N_c}\\
0&0&\cdots&0\\
\vdots&\vdots&&\vdots
\end{array}
\right)\,.
$$
以及
$$
\tilde{v}_{i}^{(i)} = \sqrt{|v_{i}^{(i)}|^2 - c^2} \, .
$$
此时规范群完全破缺；依赖于 $v_{i}^{(i)}$ 的取值，手征对称性破缺的模式相当复杂，
有多种可能。对于 $N_f > N$，有残余的 $R$ 对称性。
% [Tutor Note: 见原文 PDF 第 36 页第 2 种情形。原文标题写作 "$N_f\leq N$"，
% 但按本节开头的分类 "$N_f<N_c$ 与 $N_f\geq N_c$"，此处应为 $N_f\geq N_c$；
% 同段末尾 "For $N_f>N$, there are surviving $R$-symmetries" 中的 $N$ 也应为
% $N_c$（且按字面应是 $N_f\geq N_c$ 的情形）。两处均疑为原文把 $N_c$ 简写成 $N$
% 造成的笔误，译文照录原文。$\tilde v_i^{(i)}=\sqrt{|v_i^{(i)}|^2-c^2}$ 直接由
% 式 (77) 的对角分量 $|v_i^{(i)}|^2-|\tilde v_i^{(i)}|^2=c^2$ 得到，注意这里取
% 的是 $|\cdot|^2$ 而非 $(\cdot)^2$。]

\subsection{$N=2$ 纯规范理论中的中心荷}

我们已经看到，含中心荷 $Z$ 的 $N=2$ 超对称代数对粒子质量给出界
$M\ge \sqrt{2}|Z|$。我们还曾指出，这个界与由电荷、磁荷确定的 BPS 质量界相同。
本节我们按照 \cite{Olive-Witten} 显式计算 $Z$，以此来证明上述论断。
% [Tutor Note: 见原文 PDF 第 37 页 §2.8 开头。$M\ge\sqrt2|Z|$ 是 §2.3 讨论含中心荷
% 代数表示时得到的结论（原文 PDF 第 26--27 页，见分片 C）；本节要证明它与 §1 的
% Bogomol'nyi 界是同一个界。注意 $\sqrt2$ 因子是 $N=2$ 代数归一化的产物，本节
% 末尾（原文 PDF 第 41 页）由 (84) 之前的 $\tau_{cl}$ 表达式得到的
% $M\ge\sqrt2|Z|$ 必须与之一致，切勿中途丢掉或改成 $\sqrt2/2$。]

在超对称代数中，中心荷 $Z$ 出现在超荷 $Q^I_\alpha$ 的对易子中，而超荷本身是
$S^{I0}_\alpha$ 的空间积分（这里 $S^{I\mu}_\alpha$ 表示超流）。因此，我们首先
要由基本场算出 $S^{I0}$，然后计算它们的对易子。中心荷与这一对易子的空间积分中
的一个表面项有关；当场位形对应电荷与磁荷时，该项不为零。
  
作为热身，我们从 $N=1$ 手征超场理论开始。其拉格朗日量为
$$
{\cal L}=\int d^4 \theta\,\Phi^\dagger\Phi +\int d^2 \theta \;
{\cal W}(\Phi) + \int d^2 \bar\theta \bar{\cal W}(\Phi^\dagger)\,.
$$
定义 $y^\mu=x^\mu + i \theta \sigma^\mu \bar{\theta}$，超场可以展开为
$\Phi=A(y)+\sqrt{2}\,\theta\psi(y)
+\theta\theta F(y)$。在这个基下，$Q_{\alpha}=\partial /
\partial\theta^{\alpha}$，$\bar{Q}_{\dot{\alpha}}=-
\partial/\partial\bar{\theta}^{\dot{\alpha}} + 2i
(\theta \sigma^\mu)_{\dot{\alpha}}\,\partial/\partial y^\mu$，而场的超对称
变分 $\delta_{\epsilon}=\epsilon^\alpha Q_\alpha +
\bar{\epsilon}_{\dot{\alpha}} \bar{Q}^{\dot{\alpha}}$ 由下式给出
$$
\begin{array}{ll}
\delta A=\sqrt{2} \epsilon \psi\,,&
\delta\bar A=\sqrt{2} \bar{\epsilon}\bar{\psi}\,, \\
\delta\psi=\sqrt{2}\epsilon F + i\sqrt{2}\sigma^\mu
\bar{\epsilon} \partial_\mu A\,, &
\delta\bar{\psi}=-i\sqrt{2}\epsilon\sigma^\mu
\partial_\mu \bar{A}+\sqrt{2} \bar{F} \bar{\epsilon}\,, \\
\delta F=i\sqrt{2} \bar{\epsilon}\bar{\sigma^\mu}
\partial_\mu \psi\,, & \delta \bar{F}=i\sqrt{2}\epsilon 
\sigma^\mu \partial_\mu \bar{\psi}\,.
\end{array}
$$
利用这些变换，我们可以计算含超势 ${\cal W}$ 的项的变化为
\begin{eqnarray}
\delta ({\cal W}'F - \frac{1}{2} \frac{\partial^2 {\cal W}}{\partial
  A^2}\psi\psi) 
&=& {\cal W}''\delta AF+{\cal W}'\delta F-\frac{1}{2}{\cal W}'''\delta
A\psi\psi+{\cal W}'' \psi \delta \psi \nonumber\\
&=& {\cal W}''\sqrt{2}\epsilon\psi F+i\sqrt{2}{\cal W}'\bar{\epsilon}
\bar{\sigma}^\mu \partial_\mu \psi+\sqrt{2}{\cal W}''\psi(i\sigma^\mu 
\bar{\epsilon}\partial_\mu A + \epsilon F) \nonumber\\
&=&\partial_\mu \left( i \sqrt{2} \frac{\partial {\cal W}}{\partial A}
\bar{\epsilon} \bar{\sigma}^\mu \psi \right )\,.\nonumber
\end{eqnarray}
注意我们用了 ${\cal W}'''\delta A\psi\psi={\cal W}'''
\sqrt{2}(\epsilon \psi)\psi\psi=0$。由于 ${\cal W}'''$ 全对称，这一结论在
存在多个场 $\psi^i$ 时同样成立。于是有
$$
\delta (\int d^{2} \theta {\cal W} + \int d^{2} \bar{\theta}
\bar{{\cal W}})=i\sqrt{2}\partial_\mu\left(\frac{\partial {\cal W}}
{\partial A_j} \bar{\epsilon}\bar{\sigma}^\mu \psi^{j} -
\frac{\partial \bar{{\cal W}}}{\partial \bar{A}_j} \bar{\psi}_{j}
\bar{\sigma}^\mu \epsilon \right)\,.
$$
接下来计算动能项的变化，
$$
{\cal L}_{D} = \partial_\mu A^\dagger \partial^\mu A +
F^\dagger F-\frac{i}{2}\bar{\psi} \bar{\sigma}^\mu\partial_\mu\psi+
\frac{i}{2}\partial_\mu\bar{\psi}\bar{\sigma}^\mu\psi\,.
$$
稍微冗长的计算给出：
\begin{eqnarray}
\delta {\cal L}_{D}=&-&\frac{i}{\sqrt{2}}\partial_\mu (F\bar{\psi}
\bar{\sigma}^\mu \epsilon ) +
\frac{i}{\sqrt{2}} \partial_\mu (\epsilon\psi\partial^\mu A^\dagger - 
\epsilon \sigma^{\nu\mu}\psi \partial_\nu A^\dagger)\nonumber \\
&+& \frac{i}{\sqrt{2}}\partial_\mu (F^\dagger \bar{\epsilon}
\bar{\sigma}^\mu\psi)+\frac{i}{\sqrt{2}}\partial_\mu (\bar{\epsilon} 
\bar{\psi}\partial^\mu A-\bar{\psi}\bar{\sigma}^{\mu\nu}\bar{\epsilon}
\partial_\nu A )\,.\nonumber
\end{eqnarray}
% [Tutor Note: 见原文 PDF 第 37--38 页。原文说 "a slightly lengthy computation yields"
% 便跳过了全部中间步骤，这是本节的第一个"含糊"处：实际做法是把
% $\delta{\cal L}_D$ 中的各项分部积分，把不是全导数的部分用运动方程或
% $\sigma$ 矩阵恒等式归并，最后得到四个全导数的组合。核对要点是式 (78) 中四个
% 系数依次为 $\sqrt2$、$i\sqrt2$、$\sqrt2$、$-i\sqrt2$，其中两个 $i$ 的符号由
% $\bar\psi\bar\sigma^\mu\partial_\mu\psi$ 与 $\partial_\mu\bar\psi\bar\sigma^\mu\psi$
% 的相对符号决定；写成超流时再把 $\epsilon$、$\bar\epsilon$ 吸收进去，所以
% 实际的 $S^\rho_\alpha$ 与这里差一个 $\epsilon$ 因子（见式 (79)）。]
把超势与动能项的变化相加，并利用定义
$$
\sigma^{\mu\nu}=\frac{1}{2} (\sigma^\mu \bar{\sigma}^\nu -
\sigma^\nu \bar{\sigma}^\mu )\,,\quad
\bar{\sigma}^{\mu\nu} = \frac{1}{2} (\bar{\sigma}^\mu \sigma^\nu
- \bar{\sigma}^\nu \sigma^\mu ) \,,
$$
我们得到超流为
\begin{equation}
S^\rho_{matter}=\sqrt{2}\epsilon \sigma^\nu\bar{\sigma}^\rho \psi
\partial_\nu A^\dagger + i \sqrt{2}\frac{\partial\bar{{\cal W}}}
{\partial A^\dagger}\bar{\psi}\bar{\sigma}^\rho\epsilon
+ \sqrt{2}\bar{\psi} \bar{\sigma}^\rho\sigma^\nu\bar{\epsilon}
\partial_\nu A - i\sqrt{2}\frac{\partial {\cal W}}{\partial A} 
\bar{\epsilon} \bar{\sigma}^\rho \psi\,. 
\label{N=1scch}
\end{equation}
% [Tutor Note: 见原文 PDF 第 38 页式 (78)。底本（作者 1996 年的 LaTeX 源码）把本式
% 写在 \beq...\eeq（即 equation 环境）中，却在式中带了 \nonumber；当年的 LaTeX 2e
% 在 equation 环境中忽略 \nonumber，故原文 PDF 该式仍有编号 (78)。在载入 amsmath 的
% 现代编译下，\nonumber（=\notag）会真正取消编号，从而使 (78) 及以后所有公式的编号
% 比原文少 1。为满足“公式编号与原文逐一对应”的要求，本译本删去了这一多余的
% \nonumber（公式的数学内容一字未动），该式仍编为 (78)。]
注意为方便起见，我们把超对称变换参数 $\epsilon,\bar\epsilon$ 包含进了超流的
定义。因此按分量展开，我们实际上有：
\begin{equation}
S^\rho =\epsilon^\alpha S^\rho_\alpha + \bar\epsilon_{\dot\alpha}
\bar S^{\rho\dot\alpha}\,.
\label{S-comp}
\end{equation}
% [Tutor Note: 见原文 PDF 第 38 页式 (78)--(79)。这是本节最容易混淆的约定：
% (i) 原文把 $\epsilon,\bar\epsilon$ 吸收进 $S^\rho$，故 (79) 中 $S^\rho$ 是
% "带参数的"超流；后面式 (80) 之前的 $S^{\rho}_{gauge}$、式 (82) 以及 §2.9 中的
% $\epsilon^\alpha S^\mu_{(1)\alpha}$ 都沿用这一约定。计算超荷对易子时必须先
% 把参数提出，否则会多出 $1/\sqrt2$ 与 $i$ 因子。
% (ii) (78) 中 $\bar{\cal W}$ 与 ${\cal W}$ 的导数项分别带 $\bar\psi$ 与 $\psi$，
% 其 $i\sqrt2$ 的符号相反，这是下面 (82) 中相应项的来源。]
  
接下来考虑规范场的引入，它们由矢量多重态描述。该多重态中各场的超对称变分为
\begin{equation}
\begin{array}{lll}
\delta A^a_\mu &=& -i\bar{\epsilon}\bar{\sigma}_\mu\lambda^a +
i \bar{\lambda}^a \bar{\sigma}_\mu\epsilon \,,\\
\delta D^{a} &=&\bar{\epsilon}\bar{\sigma}^\mu D_\mu \lambda^a +
D_\mu \lambda^{a} \bar{\sigma}^\mu \epsilon \,,\\
\delta\lambda^a &=& \frac{1}{2} \sigma^{\mu\nu}\epsilon F^a_{\mu\nu} 
+ i\epsilon D^a \,,\\
\delta\bar{\lambda}^a &=&\frac{1}{2}\bar{\epsilon}\bar{\sigma}^{\nu
\mu} F^a_{\mu\nu} - i \bar{\epsilon} D^a \,.
\end{array}
\label{str1-g}
\end{equation}
此外，在存在规范相互作用时，物质场的变换需要作一些修改：
\begin{equation}
\begin{array}{lll}
\delta A &=& \sqrt{2} \epsilon \psi \,,\\
\delta \psi &=& \sqrt{2} \epsilon F + i \sqrt{2} \sigma^\mu
\bar{\epsilon} D_\mu A \,,\\
\delta\bar{\psi}&=&\sqrt{2}\bar{\epsilon} F^\dagger -i\sqrt{2}
\epsilon \sigma^\mu D_\mu A^\dagger \,,\\
\delta F &=& i\sqrt{2}\bar{\epsilon} \bar{\sigma}^\mu D_\mu \psi - 
2i T^a A \bar{\epsilon} \bar{\lambda}^a\,.
\end{array}
\label{str1-m}
\end{equation}
% [Tutor Note: 见原文 PDF 第 38 页式 (80)--(81)。这两组变分是后面全部计算的输入，
% 极易抄错：(i) $\delta A^a_\mu$ 中 $\epsilon$ 与 $\bar\epsilon$ 的位置、以及
% $-i$ 与 $+i$ 的分配；$\delta D^a$ 没有 $i$ 因子，且两项的 $\bar\sigma^\mu$ 在
% 不同的两侧。(ii) 式 (81) 第三行 $\delta\bar\psi$ 中 $\bar\epsilon$ 配 $F^\dagger$、
% $\epsilon$ 配 $D_\mu A^\dagger$，与第二行恰好"交叉"，写成同侧就会破坏超对称。
% (iii) $\delta F$ 最后一项 $-2iT^aA\bar\epsilon\bar\lambda^a$ 是为了抵消
% (69) 中 $A^\dagger T^a\lambda^a\psi$ 项里 $\delta\psi$ 的一部分而必需的，
% 系数 2 与生成元的归一化有关；若漏掉此项，超流 (82) 将不守恒。]
$\delta F$ 中的最后一项是必需的，用来抵消拉格朗日量
(\ref{fullC}) 中 $A^\dagger T^a\lambda^a\psi$ 项里 $\psi$ 变分的一部分。
描述纯矢量超场的那部分拉格朗日量可以写成
$$
{\cal L}=\frac{1}{g^2}\left(-\frac{1}{4} F^a_{\mu\nu} F^{a\mu\nu}-
\frac{i}{2}\bar\lambda^a\bar\sigma^\mu D_\mu\lambda^a+\frac{i}{2}
D_\mu \bar\lambda^a\bar\sigma^\mu\lambda^a+\frac{1}{2}D^{2} \right) + 
\frac{\theta}{32 \pi^2} F^a_{\mu\nu} \tilde{F}^{a\mu\nu} \,.
$$
若忽略 $\theta F \tilde{F}$ 项，则由拉格朗日量的这一部分得到的超流为
$$
S^\rho_{gauge}=-\frac{i}{2g^2}\left(\bar{\lambda}^a \bar{\sigma}^\rho
\sigma^{\mu\nu}\epsilon F^a_{\mu\nu}+\bar{\epsilon}\bar{\sigma}^{\mu\nu} 
\bar{\sigma}^\rho \lambda^a F^a_{\mu\nu}\right)\,.
$$
对于与物质耦合的理论，超流并不只是上述两项之和再加上 $S^\rho_{matter}$ 的
规范协变化版本。我们还预期有来自 $D^a A^\dagger T^a A$
与 $i A^\dagger T^a \lambda^a\psi+h.c.$ 这些项的贡献。把这些考虑在内，
相互作用 $N=1$ 拉格朗日量 (\ref{fullC}) 的超流为
\begin{eqnarray}
S^\rho = &-&\frac{i}{2g^2}\left(\bar{\lambda}^a \bar{\sigma}^\rho
\sigma^{\mu\nu}\epsilon +\bar{\epsilon}\bar{\sigma}^{\mu\nu} 
\bar{\sigma}^\rho \lambda^a\right) F^a_{\mu\nu} -
\left(\bar{\epsilon} \bar{\sigma}^\rho \lambda^a + \bar{\lambda}^a 
\bar{\sigma}^\rho \epsilon\right) A^\dagger T^a A
\nonumber\\
&+& \sqrt{2}\epsilon\sigma^\mu \bar{\sigma}^\rho \psi D_\mu A^\dagger
+\sqrt{2}\bar{\psi}\bar{\sigma}^\rho\sigma^\mu\bar{\epsilon}D_\mu A
+i\sqrt{2}\frac{\partial \bar{\cal W}}{\partial A^\dagger}\bar{\psi}
\bar{\sigma}^\rho\epsilon-i\sqrt{2}\frac{\partial{\cal W}}{\partial A}
\bar{\epsilon}\bar{\sigma}^\rho \psi\,.
\label{N=1sc}
\end{eqnarray}
% [Tutor Note: 见原文 PDF 第 39 页式 (82) 之前。原文那句
% "the supercurrent is not just the sum of the above, plus the
% gauge-covariantized version of $S^\rho_{matter}$" 措辞含混（直译会读成
% "不仅是上述之和再加协变化版本"），其本意是：超流并不等于"上面两项之和"
% 加上"物质超流的规范协变化"，还必须补上来自 $D^aA^\dagger T^aA$ 与
% $iA^\dagger T^a\lambda^a\psi$ 的额外项，否则超流不守恒。译文按此理解处理，
% 读者对照 (82) 与 (78) 即可看出第三、四行就是协变化后的物质项，而第二行末的
% $A^\dagger T^aA$ 项则是额外贡献。]

现在把注意力转向纯 $N=2$ 规范理论。如 §1 所讨论，它由 $N=1$ 理论
(\ref{fullC}) 令 ${\cal W}=0$ 并把手征超场 $\Phi$ 标度为 $\Phi/g$ 而得到。
此外，$\Phi$ 现在是规范群伴随表示中的一个矢量。于是把 $(\lambda,\psi)$ 重新
记为 $(\lambda_1,\lambda_2)$，我们可以容易地写下其中一个 $N=2$ 超流为
\begin{eqnarray}
g^2 S^\rho_{(1)}=&-&\frac{i}{2}\left(\bar{\lambda}^a_1\bar{\sigma}^\rho
\sigma^{\mu\nu}\epsilon + \bar{\epsilon}\bar{\sigma}^{\mu\nu}
\bar{\sigma}^\rho \lambda^a_1\right) F^a_{\mu\nu} -
\left(\bar{\epsilon}\bar{\sigma}^\rho\lambda^a_{1}+\bar{\lambda}^a_{1}
\bar{\sigma}^\rho \epsilon\right) A^\dagger T^a A 
\nonumber \\
&+&\sqrt{2}\,\epsilon\sigma^\mu\bar{\sigma}^\rho\lambda^a_{2}\, D_\mu
A^{a\dagger}+\sqrt{2}\,\bar{\lambda}^a_{2}\bar{\sigma}^\rho \sigma^\mu
\bar{\epsilon}\, D_\mu A^a \,.\nonumber 
\end{eqnarray}
$N=2$ 理论还在第二组超对称变换（参数为 $\epsilon'$）下不变，该变换由
(\ref{str1-g}) 与 (\ref{str1-m}) 作替换
$\lambda\rightarrow\psi$、$\psi\rightarrow -\lambda$ 得到。这对应于对
（$\lambda,\psi$）二重态作 $SU(2)_R$ 的一个元素变换，而 $SU(2)_R$ 正是
$N=2$ 代数的对称群。相应的守恒流于是由 $S^\rho_{(1)}$ 作替换
$\lambda_1 \rightarrow\lambda_2$、$\lambda_2\rightarrow-\lambda_1$ 得到：
\begin{eqnarray}
g^2 S^\rho_{(2)}=&-&\frac{i}{2}\left(\bar{\lambda}^a_2\bar{\sigma}^\rho
\sigma^{\mu\nu}\epsilon' + \bar{\epsilon'}\bar{\sigma}^{\mu\nu}
\bar{\sigma}^\rho \lambda^a_2 \right) F^a_{\mu\nu} -
\left(\bar{\epsilon'}\bar{\sigma}^\rho\lambda^a_{2}+\bar{\lambda}^a_{2}
\bar{\sigma}^\rho \epsilon'\right) A^\dagger T^a A 
\nonumber \\
&-&\sqrt{2}\,\epsilon'\sigma^\mu \bar{\sigma}^\rho\lambda^a_{1}\,D_\mu
A^{a\dagger}-\sqrt{2}\,\bar{\lambda}^a_{1}\bar{\sigma}^\rho \sigma^\mu
\bar{\epsilon'}\, D_\mu A^a \,.\nonumber 
\end{eqnarray}
% [Tutor Note: 见原文 PDF 第 39 页。两个超流都整体乘了 $g^2$（左端写成 $g^2S$），
% 因为下面用 $g^2S^\mu$ 与规范费米子的归一化配套；还原物理超流时别忘了除以 $g^2$。
% 第二个超流是 $SU(2)_R$ 旋转的结果：$\lambda_2\to-\lambda_1$ 使末两项整体变号，
% 同时变换参数换成 $\epsilon'$，这不是笔误。注意原文说这对应于
% "$(\lambda,\psi)$ 二重态作 $SU(2)_R$ 变换"，其具体实现是
% $(\lambda_1,\lambda_2)\to(\lambda_2,-\lambda_1)$，即 $SU(2)_R$ 的一个
% $90^\circ$ 转动（$\epsilon_{12}$ 型），读者可自行验证它保持 $N=2$ 代数不变。]
我们先集中讨论 $S^\mu_{(1)}$。利用恒等式
\begin{eqnarray}
\sigma^a \bar{\sigma}^b \sigma^c &=& \eta^{ab} \sigma^c - \eta^{ac}
\sigma^b + \eta^{bc} \sigma^a + i \epsilon^{abcd} \sigma_d \,, 
\nonumber\\
\bar{\sigma}^a \sigma^b \bar{\sigma}^c &=& \eta^{ab} \bar{\sigma}^c
- \eta^{ac} \bar{\sigma}^b + \eta^{bc} \bar{\sigma}^a - i
\epsilon^{abcd} \bar{\sigma}_d \nonumber\,,
\end{eqnarray}
以及 $\chi\sigma^\mu\bar\psi=-\bar\psi\bar\sigma^\mu\chi$，这一超流可以改写为
\begin{eqnarray}
g^2S^\mu_{(1)}=&-&\epsilon\sigma_\nu\bar{\lambda}^a_{1}(iF^{a\mu\nu}+ 
\widetilde F^{a\mu\nu})+\sqrt{2}\,\epsilon\sigma^\nu\bar{\sigma}^\mu
\lambda^a_{2}\,D_\nu A^{\dagger a}+\epsilon\sigma^\mu 
\bar{\lambda}^a_{1} A^\dagger T^a A \nonumber \\
&+& (\mbox{\it $\bar\epsilon$ dependent terms})\,.\nonumber
\end{eqnarray}
% [Tutor Note: 见原文 PDF 第 39 页。这两个 $\sigma$ 矩阵恒等式的最后一项
% $\pm i\epsilon^{abcd}$ 的符号完全取决于 $\epsilon^{0123}$ 的取值约定
% （以及 $\eta={\rm diag}(1,-1,-1,-1)$ 的号差），原文在 §1.1 的约定中固定了这一点。
% 若采用相反的 $\epsilon^{0123}$ 约定，化简后会得到 $-iF+\widetilde F$ 之类的
% 不同组合，进而使中心荷的实部/虚部归属颠倒。这是本节最容易出现"整体符号错误"
% 的地方，请务必核对原始约定。]
由上可以容易地读出分量
$S^\mu_{(1)\alpha}$（见式 (\ref{S-comp})）为
$$
g^2 S^\mu_{(1)\alpha}= \sigma_{\nu\alpha\dot\alpha}
\bar{\lambda}^{a\dot\alpha}_{1} (iF^{a\mu\nu}+\widetilde F^{a\mu\nu})
+\sqrt{2}(\sigma^\nu\bar{\sigma}^\mu\lambda^a_{2})_\alpha D_\nu 
A^{\dagger a}+\sigma^\mu_{\alpha\dot{\alpha}}
\bar{\lambda}^{a\dot{\alpha}}_{1} A^\dagger T^a A\,. 
$$
利用
$$
\bar{\lambda}^{\dot{\alpha}}_{1}=\epsilon^{\dot{\alpha}\dot{\beta}}
\bar\lambda_{1\dot\beta}=i(\sigma_y\lambda^\dagger_1)^{\dot\alpha}\,,
\qquad(\mbox{Note that}\;\bar\lambda_{\dot\beta}=\lambda^\dagger_\beta) 
$$
降低旋量指标，并对矢量的空间分量采用矢量记号，
流的 $\mu=0$ 分量取形式
\begin{eqnarray}
g^2S^0_{(1)\alpha}\,\,=&-&i(\vec{\sigma}\sigma_y\lambda^{\dagger a}_1)_\alpha
\cdot (i\vec{F}^a+\vec{\widetilde F})^a \nonumber\\
&+& \sqrt{2} \lambda^a_{2\alpha} D_0 A^{\dagger a} +
\sqrt{2}(\vec{\sigma}\cdot\vec{D}A^{\dagger a}\lambda^a_{2})_\alpha 
+i(\sigma_y\lambda^{\dagger a}_1)_\alpha A^\dagger T^a A\,, \nonumber
\end{eqnarray}
其中
$$
\vec{F}^a=F^{a 0i}\,,\quad \vec{\widetilde F}^a=\widetilde F^{a0i}\,.
$$
与前面一样，$S^0_{(2)\alpha}$ 的表达式由上式作替换
$\lambda_1\rightarrow\lambda_2$、$\lambda_2\rightarrow-\lambda_1$ 得到。
% [Tutor Note: 见原文 PDF 第 39--40 页。这一步的指标位置极易写错：
% $\sigma_{\nu\alpha\dot\alpha}\bar\lambda^{\dot\alpha}$ 与
% $(\sigma^\nu\bar\sigma^\mu\lambda)_\alpha$ 的旋量指标都作用在 $\alpha$ 上，
% 而 $\sigma^\mu_{\alpha\dot\alpha}\bar\lambda^{\dot\alpha}A^\dagger T^aA$
% 一项来自 $\epsilon\sigma^\mu\bar\lambda A^\dagger T^aA$ 中去掉 $\epsilon$。
% 另外 $\bar\lambda^{\dot\alpha}=i(\sigma_y\lambda^\dagger)^{\dot\alpha}$ 中的
% $i$ 与 $\sigma_y$ 是 $(\frac12,0)$ 与 $(0,\frac12)$ 表示互换的标准结果，
% 原文在括号里特别注明 $\bar\lambda_{\dot\beta}=\lambda^\dagger_\beta$ 供核对。]

为计算中心荷，我们关心反对易子
$$
\{Q_{(1)\alpha}\,,\,Q_{(2)\beta}\} =
\left\{\int d^3x\,S^0_{(1)\alpha}(\vec{x},0)\,,\,\int d^3y\,S^0_{(2)
\beta}(\vec{y},0)\right\} \,.
$$
正如 Olive 与 Witten \cite{Olive-Witten} 所指出的，其非零贡献来自某些测量
电荷与磁荷的边界项。为验证这一结果，我们必须在反对易子中寻找形如
$\int d^3 x\partial_i 
(A^{\dagger a} F^{aoi} + A^{\dagger a}\widetilde F^{aoi})$ 的项。
为此只需保留超流中的相关项：
\begin{eqnarray}
g^2S^0_{(1)\alpha}&=&-i(\vec{\sigma}\sigma_y\lambda^{\dagger a}_1)_\alpha
\cdot (i\vec{F}^a+\vec{\widetilde F})^a +\sqrt{2}\,(\vec{\sigma}\cdot
\vec{D}A^{\dagger a}\lambda^a_{2})_\alpha + \cdots \,,\nonumber\\
g^2S^0_{(2)\alpha}&=&-i(\vec{\sigma}\sigma_y\lambda^{\dagger a}_2)_\alpha
\cdot (i\vec{F}^a+\vec{\widetilde F})^a -\sqrt{2}\,(\vec{\sigma}\cdot
\vec{D}A^{\dagger a}\lambda^a_{1})_\alpha + \cdots\,.
\end{eqnarray}
利用它，我们得到
\begin{eqnarray}
\{Q_{(1)\alpha}\,,\,Q_{(2)\beta}\}&=&\frac{1}{g^4}\int d^3x\int d^3y\;
\big[i\sqrt{2}(\sigma^i\sigma_y)_{\alpha\gamma}\sigma^j_{\beta\lambda}
\{\lambda^{\dagger a}_{1\gamma},\lambda^b_{1\lambda}\}(iF^a_{0i}+
\widetilde F^a_{0i})D_j A^{\dagger b}     \nonumber\\
&&\qquad\qquad\quad -i\sqrt{2}(\sigma^j)_{\alpha\gamma}(\sigma^i
\sigma_y)_{\beta\lambda}\{\lambda^{\dagger a}_{2\gamma},
\lambda^b_{2\lambda}\}(iF^a_{0i}+\widetilde F^a_{0i})D_j 
A^{\dagger b}\big]      \nonumber \\
&=&\frac{1}{g^2}\int d^3x\,i\sqrt{2}\left[(\sigma^i\sigma_y
\sigma^{j\,T})_{\alpha\beta}-(\sigma^i\sigma_y\sigma^{j\,T})_{\beta\alpha}
\right](iF^a_{0i}+\widetilde F^a_{0i})D_j{A^\dagger}^{a}\,. \nonumber
\end{eqnarray}
% [Tutor Note: 见原文 PDF 第 40 页。这里原文把两处"显然"的步骤跳过了，实际用到：
% (i) 等时反对易关系 $\{\lambda^{\dagger a}_{1\gamma},\lambda^b_{1\lambda}\}
% \propto\delta^{ab}\delta_{\gamma\lambda}$（规范费米子按 $\lambda\to g\lambda$
% 归一化后，这正是第一行系数 $1/g^4$ 变成第二行 $1/g^2$ 的原因）；
% (ii) $\lambda_1$ 与 $\lambda_2$ 同属一个 $N=2$ 多重态，彼此反对易，因此两组
% 费米子不能交叉缩并，只能各自成对出现，这正是第一行两项的来历；
% (iii) $\sigma^{j\,T}$ 的出现来自交换 $\alpha\leftrightarrow\beta$ 时转置
% $\sigma^j$。最后的 $\sigma^i\sigma_y\sigma^{j\,T}$ 只有其反对称部分存活。]

方括号中涉及 $\sigma$ 矩阵的项，若利用 $\sigma^i\sigma_y=-\sigma_y\sigma^{i\,T}$，
便可化简，于是
$$
(\sigma^i\sigma_y\sigma^{j\,T})_{\alpha\beta}=-(\sigma_y\sigma^{i\,T} 
\sigma^{j\,T})_{\alpha\beta}=-[\sigma_y(\delta^{ij}-i\epsilon^{ijk}
\sigma_k^T)]_{\alpha\beta} \,.
$$
用它减去把 $\alpha$ 与 $\beta$ 互换后得到的类似式子，得到方括号内的项等于
$-2(\sigma_y)_{\alpha\beta}\delta^{ij}=2i\epsilon_{\alpha\beta}
\delta^{ij}$。因此对易子取形式
$$
\{Q_{(1)\alpha},Q_{(2)\beta}\}=-\frac{2\sqrt{2}}{g^2}
\epsilon_{\alpha\beta}\int d^3x (iF^{a0i}+\widetilde F^{a0i}) 
D_i A^{\dagger a} \,.
$$
利用磁部分的 Bianchi 恒等式与电部分的运动方程，容易证明这与下式相同
$$
\{Q_{(1)\alpha},Q_{(2)\beta}\}=-\frac{2\sqrt{2}}{g^2}
\epsilon_{\alpha \beta}\int d^3x\,\partial_i\left[(iF^{a0i}+
\widetilde F^{a0i}) A^{\dagger a}\right] \,.
$$
% [Tutor Note: 见原文 PDF 第 40--41 页。两点提示：(i) 由
% $-2(\sigma_y)_{\alpha\beta}\delta^{ij}$ 写成 $2i\epsilon_{\alpha\beta}\delta^{ij}$
% 用了 $(\sigma_y)_{\alpha\beta}=-i\epsilon_{\alpha\beta}$ 的约定（即
% $\epsilon_{12}=+1$、$\sigma_y=\begin{pmatrix}0&-i\\ i&0\end{pmatrix}$）；
% 若把 $\epsilon_{\alpha\beta}$ 的符号约定弄反，中心荷的实部与虚部就会对调。
% (ii) 从 $F^{a0i}D_iA^{\dagger a}$ 到全导数 $\partial_i(F^{a0i}A^{\dagger a})$
% 需要：对磁部分用 Bianchi 恒等式 $D_i\widetilde F^{a0i}=0$，对电部分在源外区域
% 用运动方程 $D_iF^{a0i}=0$。原文只说 "one can easily show"，这一步实质上是把
% $\partial_i(F A)$ 展开后丢掉 $D_iF$ 项；它正是 §1 中 Bogomol'nyi 表面项论证
% 的超对称版本。]
类似地，
$$
\{\bar Q_{(1)\dot{\alpha}},\bar Q_{(2)\dot{\beta}}\}=-\frac
{2\sqrt{2}}{g^2} \epsilon_{\dot{\alpha}\dot{\beta}} \int d^3 x\,
\partial_i \left[(-iF^{a0i}+\widetilde F^{a0i}) A^a \right]\,.
$$
于是，在超流的对易子中，我们重新得到了全导数，它们正是电荷密度与磁荷密度：
$$
Q_{ele}=-\frac{1}{ag}\int d^3x \partial_i(F^{a0i} A^a)= gn_e\,,\qquad 
Q_{mag}=-\frac{1}{ag}\int d^3x \partial_i(\widetilde F^{a0i} A^a)=
\frac{4\pi}{g}n_m \,.
$$
% [Tutor Note: 见原文 PDF 第 41 页电荷与磁荷的表达式。这是全节最易出归一化错误
% 的地方：(i) $Q_{ele}$ 与 $Q_{mag}$ 的分母都是 $ag$，但右端分别是 $gn_e$ 与
% $\frac{4\pi}{g}n_m$；磁荷的 $4\pi/g$ 来自 't Hooft 张量 $\widetilde F$ 的归一化
% （见 §1.4 关于电荷量子化的讨论，原文 PDF 第 7 页起）。(ii) $a$ 是 Higgs 真空中
% $A$ 的取值，其归一化取决于生成元的选取（原文 $T^a$ 取基本表示，
% ${\rm Tr}(T^aT^b)=\frac12\delta^{ab}$）；若换成 $A=\frac a2\sigma_3$ 的写法，
% $Q_{ele}$ 的系数就会差因子 2。(iii) 两个公式都隐含了 $a$ 与 $g$ 的乘积
% $ag$，这正是狄拉克量子化条件 $eg=2\pi$ 型结构的体现。]
这里 $a$ 是 $A$ 在 Higgs 真空中的取值。上面用到的电荷量子化条件对应于积分的
基本电荷，正如 $SU(2)$ 破缺到 $U(1)$ 且场处于伴随表示的情形。现在与
(\ref{Susy-C'}) 比较，可以容易地读出 $N=2$ 中心荷为 $Z=-ia(n_e+ (4\pi
i/g^2)n_m)$。回忆 $Z$ 的相位是依赖于约定的，我们主要关心它的模长。
% [Tutor Note: 见原文 PDF 第 41 页。此处 $Z=-ia(n_e+(4\pi i/g^2)n_m)$ 与下面
% 含 $\theta$ 的 $Z=a(n_e+\tau_{cl}n_m)$ 相差一个整体因子 $-i$（原文自己也说
% "ignoring an overall factor of $-i$"）。两种写法不可混用：$|Z|$ 不受整体相位
% 影响，但 BPS 等号条件 $M=\sqrt2|Z|$ 中 $Z$ 的相位与超对称投影（哪一个超荷
% 被保留）有关，写有效作用量时一般用后面 $Z=an_e+a_Dn_m$ 的全纯形式。
% 另外此处的比较对象是式 (52)（原文 PDF 第 22 页附近，分片 C 的范围内），
% 原文说 "comparing with (52)" 即指 §2.2 中中心扩张的代数形式。]
为找出 $\theta$ 参数的影响，我们可以修改上述计算，把 $\theta F\widetilde F$
的贡献加进去，或者直接利用 §1.10 中描述的 Witten 效应。如那里所学到的，
$\theta$ 参数的作用是把电荷移动为
$Q_{ele}=gn_e+(\theta g^2/8\pi^2)Q_{mag}$。利用这一点（并忽略一个整体的
$-i$ 因子），我们得到中心荷为
$$
Z= a\,(n_e + \tau_{cl} n_m)\,,\quad\rm{where},\quad 
\tau_{cl} = \frac{\theta}{2\pi} + \frac{4\pi i}{g^2}\,.
$$
% [Tutor Note: 见原文 PDF 第 41 页。(i) 原文 "the Witten effect described in
% section (1.10)" 指的是 §1.10 “$\theta$ 参数与磁单极电荷”（原文 PDF 第 20 页起），
% 该处正是推导电荷平移 $q=ne+\theta e^2g/8\pi^2$ 的地方；原文编号无误。
% （本讲义 §1 共 10 个小节，§1.10 即 $\theta$ 参数一节。）
% (ii) 复耦合 $\tau_{cl}=\theta/2\pi+4\pi i/g^2$ 的下标 $cl$ 表示这是经典值
% （未含量子修正）；它与 (72) 第一行的 $\tau$ 一致，区别只在本书后文会用
% 量子修正后的 $\tau$。
% (iii) $Q_{ele}$ 的移动量 $(\theta g^2/8\pi^2)Q_{mag}$ 与 $Q_{mag}=4\pi n_m/g$
% 合起来给出 $Q_{ele}=g(n_e+\frac{\theta}{2\pi}n_m)$，这正是 $\tau_{cl}$ 的
% 实部来源；核对时注意 $\theta$ 的周期 $2\pi$。]
利用超对称代数，这意味着质量界
$$
M \geq \sqrt{2}\, |Z|= {\sqrt 2}\,|\,a\,(n_e + \tau_{cl} n_m )\,|\,,
$$
即 BPS 界。至此我们考虑的是微观 $N=2$ 作用量 (\ref{N=2})。在由预势 ${\cal F}$
指定的有效作用量 (\ref{N=2general}) 层面上，中心荷取形式
$$
Z = an_e + a_D n_m \,,
$$
其中 $a_D=\partial{\cal F}/\partial a$。这个公式可以如下理解：正如我们将在
下一节看到的，理论 (\ref{N=2general}) 有一种对偶描述，其中基本对象是磁单极
而不是电荷。在这一描述中，$n_e$ 与 $n_m$ 互换，$a$ 也换成 $a_D$。因此在对偶
理论中，容易推断出磁单极对 BPS 界的贡献为 $a_Dn_m$。正如我们将在下一节看到的，
完整的对偶群是 $SL(2,Z)$，在它之下 $a$ 与 $a_D$ 作为二重态变换。把这些事实
结合起来，就得到上面给出的中心扩张公式。类似的构造适用于带伴随物质的任意群
$SU(N)$。
% [Tutor Note: 见原文 PDF 第 41 页末段。(i) 原文用对偶论证而不是直接计算来"导出"
% $Z=an_e+a_Dn_m$，这属于典型的"含糊点拨"处：严格证明需要 §3 的 $SL(2,Z)$
% 变换性质（原文 PDF 第 43 页起，分片 E），这里只需记住 $(a,a_D)$ 在
% $SL(2,Z)$ 下按二重态变换、且 $a_D=\partial{\cal F}/\partial a$。
% (ii) 这里的 $a$ 是标量场 $A$ 的真空期望值（在 $SU(2)$ 中 $a$ 是
% Cartan 子代数分量的值），$a_D$ 是"对偶"标量；$Z$ 用 $a$ 与 $a_D$ 表达是
% Seiberg--Witten 理论的核心公式之一。
% (iii) 注意 $\tau$ 与预势的关系：$\tau=\partial^2{\cal F}/\partial a^2$，
% 因此 $a_D=\partial{\cal F}/\partial a$ 与 $\tau$ 是同一预势的一阶、二阶导数，
% 后文 §3.6--3.7（原文 PDF 第 54、57 页）的单值变换与 $U(1)$ 的 $\beta$ 函数
% 都建立在这个关系上。]

\subsection{含物质的 $N=2$ 规范理论中的中心荷}

当把基本表示的物质加入上面描述的 $N=2$ 纯规范理论时，有一处重要的变化：
中心荷也会收到来自物质场质量的贡献。在 $N=2$ 超对称中，物质场是超多重态的
一部分。用 $N=1$ 记号描述超多重态最方便的方式是：它由一个手征超场 $Q$ 与一个
反手征超场 $\widetilde Q^\dagger$ 组成。$Q$（$\widetilde Q$）包含分量
$q,\psi_q$（$\widetilde q,\psi_{\widetilde q}$），并在规范群 $SU(N)$ 的
$N$（$\bar{N}$）表示下变换。在 $N=2$ 代数的 $SU(2)_{R}$ 对称性下，
$q,\widetilde q^\dagger$ 构成二重态，而 $\psi_q$ 与
$\psi^\dagger_{\tilde q}$ 是单态。由于超多重态中各场的自旋 $\leq 1/2$，
它们属于 $N=2$ 代数的短表示，因此必须满足关系 $M=\sqrt{2} Z$。上一节导出的
$Z$ 在存在裸超多重态质量时不满足这一关系，这表明中心荷必须收到来自超多重态的
进一步贡献。下面我们就来计算这一贡献。
% [Tutor Note: 见原文 PDF 第 42 页 §2.9 开头。原文写短多重态 "must satisfy
% the relation $M=\sqrt{2}Z$"，严格说应为 $M=\sqrt{2}|Z|$：$Z$ 含相位（§2.8 末尾
% 已指出其相位约定任意），短（BPS）多重态的质量只由 $|Z|$ 决定；原文此处的
% $Z$ 被默认为正实数（在适当相位约定下）。译文照录原文。另注意此处逻辑：
% "$Z$ 不满足该关系" 是推出"中心荷还缺一块"的依据，而不是说前面的计算有错。]

除了手征场 $Q$ 与 $\widetilde Q$ 的标准动能项与规范耦合之外，含物质的 $N=2$
拉格朗日量还包含由下式给出的 $N=1$ 超势
$$
\sum_{i=1}^{N_f}\sqrt{2}\widetilde{Q}_i\Phi Q_i+\sum_{i=1}^{N_f} m_i 
\widetilde{Q}_i Q_i + h.c.
$$
这里指标 $i$ 遍历 $N_f$ 个夸克味，$\Phi$ 是矢量多重态在伴随表示中的手征超场。
第一项通过 $N=2$ 超对称性与物质场的规范耦合相联系，第二项则是 $N=2$ 不变的质量项。
% [Tutor Note: 见原文 PDF 第 42 页超势。请注意两项的相对系数：$\sqrt2\widetilde QQ\Phi$
% 与 $m\widetilde QQ$ 必须同时出现且系数固定，否则 $N=2$ 会破缺为 $N=1$；
% 常见错误是只写质量项而忘记 $\sqrt2\widetilde Q\Phi Q$，或把 $\sqrt2$ 写成 2。
% 这与 §2.7 式 (75) 是同一个超势，只是这里写出了味指标求和并明确 $Q_i$ 在
% $N$、$\widetilde Q_i$ 在 $\bar N$ 表示。]
当所有质量都相等时，理论还具有 $SU(N_f)$ 味对称性；当质量不相等时，它破缺到
更小的子群。若所有质量互不相等，它破缺到 $U(1)^{N_f}$。

$Q$ 与 $\widetilde Q$ 对超流的贡献可以用 (\ref{N=1scch}) 容易地算出，其中超势
${\cal W}$ 如上所给。于是作为对 $S^\mu_{(1)\alpha}$ 的贡献，我们得到（略去色
指标与味指标）：
\begin{eqnarray}
\epsilon^\alpha S^\mu_{(1)\alpha}=\cdots&+&\sqrt{2}\,\epsilon\sigma^\nu  
\bar{\sigma}^\mu\psi_q D_\nu q^\dagger+i\sqrt{2}\,m\,
\widetilde{q}^\dagger \bar{\psi}_q\bar{\sigma}^\mu \epsilon\nonumber\\
&+& \sqrt{2}\,\epsilon\sigma^\nu\bar{\sigma}^\mu\psi_{\widetilde{q}}
D_\nu {\widetilde{q}^\dagger} + i \sqrt{2}\, m\, q^\dagger
\bar{\psi}_{\widetilde{q}} \bar{\sigma}^\mu \epsilon +\cdots\,.\nonumber
\end{eqnarray}
为得到第二个超流中相应的项，我们必须作替换
$q\rightarrow\widetilde q^\dagger$、$\widetilde q^\dagger\rightarrow -q$，
\begin{eqnarray}
\epsilon^\alpha S^\mu_{(2)\alpha}=\cdots&+&\sqrt{2}\,\epsilon\sigma^\nu  
\bar{\sigma}^\mu\psi_q D_\nu \widetilde q - i\sqrt{2}\,m\,
q \bar{\psi}_q\bar{\sigma}^\mu \epsilon\nonumber\\
&-& \sqrt{2}\,\epsilon\sigma^\nu\bar{\sigma}^\mu\psi_{\widetilde{q}}
D_\nu q + i \sqrt{2}\, m\, \widetilde q \bar{\psi}_{\widetilde{q}} 
\bar{\sigma}^\mu \epsilon +\cdots\nonumber\,.
\end{eqnarray}
% [Tutor Note: 见原文 PDF 第 42 页这两组式子。要点：(i) 替换规则
% $q\to\widetilde q^\dagger$、$\widetilde q^\dagger\to-q$ 只作用于标量
% （$q,\widetilde q^\dagger$ 构成 $SU(2)_R$ 二重态），费米子
% $\psi_q,\psi_{\widetilde q}$ 是单态因而不变；这与 §2.8 中
% $\lambda_1\to\lambda_2,\lambda_2\to-\lambda_1$ 是同一个 $SU(2)_R$ 旋转。
% (ii) $m$ 项前的 $i\sqrt{2}$ 因子来自式 (78) 中
% $i\sqrt{2}\,\partial\bar{\cal W}/\partial A^\dagger\,\bar\psi\bar\sigma^\rho\epsilon$
% 一项；因为 $\partial\bar{\cal W}/\partial\bar q = m\widetilde q$ 型，
% 代入后即得这些项。(iii) 第二组式子中 $m$ 项的整体符号翻转，正是因为
% $\widetilde q^\dagger\to-q$；若符号写错，反对易子中的附加项
% $\sum_i m_i S_i$ 就会抵消而不是相加。]
利用标准的正则对易关系，并按照与上一节类似的步骤，可以计算这些额外项对
超荷反对易子的贡献。附加项的形式为
$2i\epsilon_{\alpha\beta}\,\sum_im_i S_i$，其中
$$
S_i=\int d^3x\left(D_0q_i^\dagger q_i + q_iD_0q_i^\dagger
-\frac{i}{2}\psi^\dagger_{qi}\psi_{qi} + \frac{i}{2}\psi_{qi}
\psi^\dagger_{qi} -(q\rightarrow \widetilde q, \psi_q\rightarrow
\psi_{\widetilde q})\right)\,,
$$   
这里指标 $i$ 不求和。显然，这是一个与整体 $U(1)$ 对称性相应的守恒荷，在该
对称性下 $Q_i$ 与 $\widetilde Q_i$ 分别带电荷 $+1$ 与 $-1$。这些正是破缺的
味群的 $U(1)$ 因子。
% [Tutor Note: 见原文 PDF 第 42--43 页 $S_i$ 的定义。这是标准的 Noether 荷：
% 第一项 $D_0q^\dagger q+qD_0q^\dagger$ 是标量场的"动量"部分（对称化以保证
% 厄米性），中间两项是费米子双线性的自旋贡献，括号内最后的替换
% $(q\to\widetilde q,\ \psi_q\to\psi_{\widetilde q})$ 表示对 $\widetilde q$
% 取相反的 $U(1)$ 荷（即整项变号）。原文说 "$i$ is not summed over"，
% 即 $S_i$ 对应第 $i$ 个味，而非求和。注意 $\psi^\dagger_{qi}\psi_{qi}$ 与
% $\psi_{qi}\psi^\dagger_{qi}$ 的次序在两项中相反，这是费米子算符次序所致，
% 保持原文即可。]
把这一附加项考虑在内，中心荷的公式取形式
\begin{equation}
Z = n_e a + n_m a_D + \sum_{i} \;\frac{1}{\sqrt{2}} \; m_i S_i \,.
\label{bps-m}
\end{equation}
这一公式在含夸克超多重态的 $N=2$ $SU(2)$ 规范理论的 Seiberg--Witten 分析中
是至关重要的。
% [Tutor Note: 见原文 PDF 第 43 页式 (84)。这是本片的最后结果，也是 §2 的核心公式
% 之一。系数 $1/\sqrt2$ 的来源值得记住：对 BPS 态 $M=\sqrt2|Z|$，而超多重态的
% 质量项直接进入 $M$，故质量 $m_i$ 必须以 $m_i/\sqrt2$ 的形式出现在 $Z$ 中，
% 再由 $S_i=\pm1$（对 $U(1)$ 荷为 $\pm1$ 的态）给出 $M=m_i$。若把 $1/\sqrt2$
% 写成 $\sqrt2$，BPS 谱就会整体差因子 2。另外式 (84) 是 §3、§4 中确定
% Seiberg--Witten 曲线质量参数的基础：裸质量 $m_i$ 正是通过这一项与
% $a,a_D$ 一起进入中心荷，从而决定 BPS 谱的谱线。]
